\section{问题重述}

\subsection{问题背景}

随着 AI 训练与推理规模持续扩大，多区域数据中心需要同时处理异构算力调度、网络时延和能源协同问题。已有工作分别讨论了数据中心负荷的时空可转移特性、集群灵活边界下的调度、多数据中心时空协同以及多元灵活性协同运行\cite{chen_idc_2022,zhang_flex_dro_2024,wen_vpp_2024,zhou_multiflex_2024}。跨地域计算中的任务调度已有系统综述和大规模机器学习训练实践；碳感知工作负载迁移研究也同时讨论了负载弹性的利用及其适用边界\cite{wu_survey_2025,choudhury_mast_2024,hanafy_carbonscaler_2023,sukprasert_limitations_2024}。本题中不同区域在算力、PUE、电价、新能源及储能条件上存在显著差异，任务迁移会改变区域用电负荷，储能与购售电又会影响调度的经济性。因此，本文依次研究需求预测、任务调度、储能调度及算电联合优化四个子问题。

\begin{figure}[!htbp]
\centering
\includegraphics[width=0.6\textwidth]{\ProjectRoot/modules/10_restatement/figure/背景图.jpg}
\caption{AI云数据中心跨区域调度与能源协同优化场景示意}
\label{fig:q-background}
\end{figure}

\subsection{任务要求}

\begin{enumerate}[label=\textbf{问题\arabic*.},leftmargin=4.2em,itemsep=0.45ex,topsep=0.35ex]
  \item 预测未来 24 h GPU 需求，并对第 2376--2399 小时实际任务进行基础调度，给出甘特图和各区域 GPU 利用率。

  \item 在考虑电价、碳强度与新能源的条件下优化任务执行区域和开工时刻，比较运行成本、碳排放、网络时延及新能源利用率。

  \item 固定计算负荷，优化新能源分配、电网购售电和储能充放电，分析储能对成本、碳排放、峰值净购电和电网波动的影响。

  \item 联合优化任务迁移、开工时刻、储能和购售电，在六项指标间进行权衡，并比较不同碳约束、电价机制和新能源波动下的策略变化。
\end{enumerate}